\documentclass{article}
\usepackage{amsmath, amssymb, amsthm}
\usepackage{geometry}
\geometry{a4paper, margin=1in}
\usepackage{fancyhdr}
\usepackage{hyperref}
\usepackage[russian]{babel}
\usepackage[utf8]{inputenc}
\usepackage{listings}
\usepackage{xcolor}

\pagestyle{fancy}
\fancyhf{}
\fancyhead[L]{Algorithms}
\fancyhead[R]{\thepage}

\title{Algorithms 1 \\ Homework 5}
\author{Your Name}
\date{\today}

\lstset{
    language=Python,
    basicstyle=\ttfamily\small,
    keywordstyle=\color{blue},
    stringstyle=\color{red},
    commentstyle=\color{gray},
    backgroundcolor=\color{lightgray!20},
    frame=single,
    rulecolor=\color{black},
    numbers=left,
    numberstyle=\tiny\color{gray},
    stepnumber=1,
    numbersep=10pt,
    showspaces=false,
    showstringspaces=false,
    breaklines=true,
    breakatwhitespace=true,
    tabsize=4,
    captionpos=b
}

\begin{document}

\maketitle

\section{Статистики, битовые операции, нижние оценки}
\begin{enumerate}
    \item Для решения данной задачи можно хранить
          две переменные: $cur_elem$, $running_count$,
          которые будут хранить текущий элемент и количество
          его вхождений. Алгоритм итерируется по массиву и
          на каждом шаге проверяет, равен ли текущий элемент
          $cur_elem$. Если равен, то увеличиваем $running_count$,
          иначе уменьшаем $running_count$. Если $running_count = 0$,
          то присваиваем $cur_elem$ текущему элементу. В конце
          алгоритма искомый элемент будет храниться в $cur_elem$.
          Этот алгоритм работает за $O(n)$ времени и $O(1)$ памяти.
          \begin{lstlisting}
def majority_element(arr):
    cur_elem = None
    running_count = 0
    for elem in arr:
        if running_count == 0:
            cur_elem = elem

        running_count += 1 if cur_elem == elem else -1

    return cur_elem
    \end{lstlisting}

    \item Асимптотика алгоритма поиска $k$-той
          порядковой статистики при разбиении массива
          на подмассивы из $p$ элементов:
          \[
              T(n) = T(\frac{n}{p}) + Cn + T(\frac{3p - 1}{4p} n)
          \]
          где $C$ --- константа, $n$ --- размер массива,
          Если предположить, что сортировка $p$ элементов
          занимает $O(1)$ времени.
          Асимптотика будет представлять собой дерево
          рекурсии глубины от $\log_p n$ до
          $\log_{\frac{4p}{3p - 1}} n$, асимптотика
          алгоритма на $i$-ом уровне рекурсии
          (считая с вершины дерева) будет равна
          $Cn(\frac{3p + 3}{4p})^i$, таким образом
          Асимптотика алгоритма будет равна
          \[
              T(n) = O(\sum_{i=1}^{\lceil \log_{\frac{4p}{3p - 1}}{n} \rceil} n(\frac{3p + 3}{4p})^i)
          \]
          Для $\frac{3p + 3}{4p} < 1$ сумма будет
          ограничена сверху константой, тогда асимптотика
          будет равна $O(n)$.

          Подставим $p = 3$:
          \[
              T(n) = O(\sum_{i=1}^{\lceil \log_{\frac{3}{2}}{n} \rceil} n1^i) = O(n \log n)
          \]

          Подставим $p = 7$:
          \[
              T(n) = O(\sum_{i=1}^{\lceil \log_{\frac{7}{5}}{n} \rceil} n(\frac{6}{7})^i) = O(n)
          \]

    \item Этот процесс эквивалетен поиску наибольшего
          общего делителя алгоритмом евклида. Рассмотрим
          пару чисел $a, b$, без ограничения общности
          $a > b$. Тогда $a = bq + r$, после применения
          данной операции $q$ раз пара чисел станет равна
          $b, a \mod b$. Теперь докажем, что этот процесс
          сводится к применению алгоритма евклида попарно
          ко всем элементам. Заметим, что
          $\gcd(a, b) = \gcd(a - b, b)$, тогда процесс
          можно представить как применение этого процесса
          сначала к $a_1, a_2$ до нахождения их $\gcd$,
          затем к $a_1, a_3$ и так далее. Таким образом,
          процесссводится к поиску $\gcd$ попарно ко всем э
          лементам, и он гарантированно сойдется
          к $\gcd$ всех элементов массива.

    \item Пусть $x$ и $y$ - битовые длины данных чисел.
          Тогда асимптотическая сложность этого алгоритма:
          \[
              T(x, y) = T(x - 1, y) + O(x + y)
          \]
          Так как рекурсивный вызов происходит с уменьшением
          длины одного из чисел на 1, после чего с обоими
          числами проводятся линейные операции: сложение,
          вычитание, сравнение, умножение на 2 чисел
          $q$ и $r$, линейно зависящих от $x$. Тогда
          асимптотика данного алгоритма:
          \[
              T(x, y) = O(x(x + y))
          \]

    \item Для примера будем использовать алгоритм
          быстрого возведения в степень:
          \begin{lstlisting}
def power(a, n, m): # n is bit representation
    res = 1
    a = a % m
    for i in range(len(n)):
        if n[i] == "1":
            res = (res * a) % m
        a = (a * a) % m

    return res
    \end{lstlisting}
          Возьмём $a, n, m$ - длины
          битовых представлений чисел $a, n, m$.
          Тогда асимптотика данного алгоритма:
          \[
              T(a, n, m) = = O(na(a + m))
          \]
          Алгоритм проходит по всем битам числа $n$,
          на каждой итерации в худшем случае производит
          умножение числа $a$ на $res$ по модулю $m$,
          и умножение на $a$ по модулю $m$, умножение
          происходит за $O(a^2)$ деление по модулю за
          $O(a(a + m))$.

    \item Оптимальный алгоритм делит исходную
          кучу пополам, проверяет каждую счётчиком за константу,
          потом рекурсивно проводит ту же операцию для
          кучи, где счётчик показал радиацию. Такой алгоритм
          будет работать за 
          \[ 
          O(\lceil \log n \rceil)
          \]
          Если считать время
          измерения счётчиком за константу.
          Докажем, что меньшее количество операций сделать
          невозможно. Пусть одно применение счётчика присваивает
          значение $1$ ко всем камням кучи, где найдена радиация,
          и $0$ камням, где радиации нет. Тогда по завершению
          процесса каждому камню будет присвоена уникальная
          битовая маска длины $\lceil \log n \rceil$.
          Предположим, что можно было бы провести меньше
          операций, например, $\lceil \log n \rceil - 1$.
          Тогда по принципу Дирихле найдутся два камня с одинаковыми
          масками, что противоречит тому, что каждому камню
          присвоена уникальная маска. Таким образом,
          оптимальное количество операций равно $\lceil \log n \rceil$.
\end{enumerate}
\end{document}